\documentclass[12pt]{article}
\usepackage{amsmath, amssymb, geometry, xcolor, mdframed, enumerate}
\geometry{letterpaper, margin=1in}

\title{\textbf{Homework 2 Logic - Quantified Statements and Chapter Review}}
\author{Tashfeen Omran}
\date{September 2026}

\begin{document}
\maketitle
\tableofcontents
\newpage

\section{Part 1: Comprehensive Introduction, Context, and Prerequisites}

\subsection{Core Concepts}
This chapter delves into the formal machinery of logical statements, focusing extensively on \textbf{quantifiers}, \textbf{logical equivalence}, and \textbf{implications}. While basic propositional logic uses fixed statements (e.g., $P$: ``The integer 3 is odd''), predicate logic uses open sentences (e.g., $P(x): x > 3$) whose truth value depends on a variable within a specified domain. 
To convert open sentences into statements with definitive truth values, we use quantifiers:
\begin{itemize}
    \item \textbf{Universal Quantifier} ($\forall$): Read as ``for all'' or ``for every.'' A statement $\forall x \in S, P(x)$ is true if and only if $P(x)$ is true for every single element in the domain $S$.
    \item \textbf{Existential Quantifier} ($\exists$): Read as ``there exists'' or ``for some.'' A statement $\exists x \in S, P(x)$ is true if there is at least one element in $S$ for which $P(x)$ is true.
\end{itemize}

\subsection{Intuition and Derivation}
The intuition behind quantifiers stems from generalizations of the \textit{AND} ($\land$) and \textit{OR} ($\lor$) operators. If a domain $S$ has a finite number of elements $\{a_1, a_2, \dots, a_n\}$, then:
\[ \forall x \in S, P(x) \quad \text{is logically equivalent to} \quad P(a_1) \land P(a_2) \land \dots \land P(a_n) \]
\[ \exists x \in S, P(x) \quad \text{is logically equivalent to} \quad P(a_1) \lor P(a_2) \lor \dots \lor P(a_n) \]
Understanding this connection immediately explains how to negate quantified statements. Just as De Morgan's Laws state $\sim(P \land Q) \equiv \sim P \lor \sim Q$, negating a universal statement flips the \textit{AND}s to \textit{OR}s, meaning $\forall$ becomes $\exists$.

\subsection{Historical Context and Motivation}
The formalization of quantifiers occurred in the late 19th and early 20th centuries, primarily through the works of Gottlob Frege, Charles Sanders Peirce, and Bertrand Russell. Mathematicians were struggling with the ``crisis of foundations''---a period where intuitive definitions in calculus (like limits and continuity) led to paradoxes and contradictions. They needed a rigorous, symbolic language to define concepts unambiguously. The invention of the existential and universal quantifiers allowed mathematicians to write the precise $\epsilon$-$\delta$ definitions of limits, firmly grounding the calculus you study today.

\subsection{Key Formulas and Logical Equivalences}
\begin{enumerate}
    \item \textbf{De Morgan's Laws for Quantifiers}:
    \begin{align*}
        \sim(\forall x \in S, P(x)) &\equiv \exists x \in S, \sim P(x) \\
        \sim(\exists x \in S, P(x)) &\equiv \forall x \in S, \sim P(x)
    \end{align*}
    \item \textbf{Negating an Implication}:
    \[ \sim(P \implies Q) \equiv P \land \sim Q \]
    \item \textbf{Contrapositive of an Implication} (logically equivalent to the original):
    \[ P \implies Q \equiv \sim Q \implies \sim P \]
    \item \textbf{Biconditional Characterization}:
    \[ P \iff Q \equiv (P \implies Q) \land (Q \implies P) \]
\end{enumerate}

\subsection{Prerequisites}
To master this section, you must have a solid grasp of:
\begin{itemize}
    \item \textbf{Set Theory Notation}: Understanding domains ($\mathbb{R}$, $\mathbb{Z}$, $\mathbb{Q}$, $\mathbb{N}$) and set relations ($\in$, $\subset$, $\cap$, $\cup$).
    \item \textbf{Basic Algebra}: Solving simple polynomial equations to find ``witnesses'' for existential claims.
    \item \textbf{Truth Tables}: Being able to mechanically evaluate the truth values of compound statements using conjunction, disjunction, and implication.
\end{itemize}

\newpage
\section{Part 2: Detailed Homework Solutions}
\textit{Note: The assigned scope is Section 2.10 Exercises (2.67--2.81), Section 2.11 Exercises (2.82--2.86), and the Chapter 2 Supplemental Exercises (2.87--2.107).}

\subsection{Section 2.10 Exercises}

\textbf{Exercise 2.67}
Let $S$ denote the set of odd integers and let $P(x): x^2 + 1 \text{ is even}$ and $Q(x): x^2 \text{ is even}$. State $\forall x \in S, P(x)$ and $\exists x \in S, Q(x)$ in words.
\begin{mdframed}
\textbf{Solution:} 
The statement $\forall x \in S, P(x)$ translates to: ``For every odd integer $x$, the quantity $x^2+1$ is an even integer.'' 
The statement $\exists x \in S, Q(x)$ translates to: ``There exists an odd integer $x$ such that $x^2$ is an even integer.''
\end{mdframed}

\textbf{Exercise 2.68}
Define an open sentence $R(x)$ over some domain $S$ and then state $\forall x \in S, R(x)$ and $\exists x \in S, R(x)$ in words.
\begin{mdframed}
\textbf{Solution:}
Let $S = \mathbb{R}$ (the set of real numbers). Let $R(x): x^2 \ge 0$.
$\forall x \in \mathbb{R}, R(x)$: ``For every real number $x$, its square is greater than or equal to zero.''
$\exists x \in \mathbb{R}, R(x)$: ``There exists a real number $x$ such that its square is greater than or equal to zero.''
\end{mdframed}

\textbf{Exercise 2.69}
State the negations of the following quantified statements, where all sets are subsets of some universal set $U$.
(a) For every set $A$, $A \cap \bar{A} = \emptyset$. 
(b) There exists a set $A$ such that $\bar{A} \subseteq A$.
\begin{mdframed}
\textbf{Solution:}
Using De Morgan's Laws for quantifiers ($\forall \to \exists$, $\exists \to \forall$):
(a) \textbf{Final Answer:} There exists a set $A$ such that $A \cap \bar{A} \neq \emptyset$.
(b) \textbf{Final Answer:} For every set $A$, $\bar{A} \not\subseteq A$.
\end{mdframed}

\textbf{Exercise 2.70}
State the negations of the following quantified statements:
(a) For every rational number $r$, the number $1/r$ is rational.
(b) There exists a rational number $r$ such that $r^2 = 2$.
\begin{mdframed}
\textbf{Solution:}
(a) \textbf{Final Answer:} There exists a rational number $r$ such that $1/r$ is not rational. (Note: $r=0$ is a valid witness making this negated statement true).
(b) \textbf{Final Answer:} For every rational number $r$, $r^2 \neq 2$.
\end{mdframed}

\textbf{Exercise 2.71}
Let $P(n): (5n - 6)/3 \text{ is an integer}$. Be an open sentence over the domain $\mathbb{Z}$. Determine whether the following are true:
(a) $\forall n \in \mathbb{Z}, P(n)$. \quad (b) $\exists n \in \mathbb{Z}, P(n)$.
\begin{mdframed}
\textbf{Solution:}
(a) To test $\forall n \in \mathbb{Z}$, try an arbitrary integer, say $n=0$. Then $(5(0)-6)/3 = -6/3 = -2 \in \mathbb{Z}$. Try $n=1$: $(5(1)-6)/3 = -1/3 \notin \mathbb{Z}$. Since it fails for $n=1$, the universal statement is false.
(b) To test $\exists n \in \mathbb{Z}$, we need only one witness. We already found $n=0$ works. 
\textbf{Final Answer:} (a) False. (b) True.
\end{mdframed}

\textbf{Exercise 2.72}
Determine the truth value of each of the following statements.
(a) $\exists x \in \mathbb{R}, x^2 - x = 0$
(b) $\forall n \in \mathbb{N}, n + 1 \ge 2$
(c) $\forall x \in \mathbb{R}, \sqrt{x^2} = x$
(d) $\exists x \in \mathbb{Q}, 3x^2 - 27 = 0$
(e) $\exists x \in \mathbb{R}, \exists y \in \mathbb{R}, x + y + 3 = 8$
(f) $\forall x, y \in \mathbb{R}, x + y + 3 = 8$
(g) $\exists x, y \in \mathbb{R}, x^2 + y^2 = 9$
(h) $\forall x \in \mathbb{R}, \forall y \in \mathbb{R}, x^2 + y^2 = 9$
\begin{mdframed}
\textbf{Solution:}
(a) True (witnesses: $x=0, x=1$).
(b) True (smallest natural number is $1$, so $1+1 \ge 2 \implies 2 \ge 2$, which is true).
(c) False (counterexample: $x=-1$. $\sqrt{(-1)^2} = 1 \neq -1$).
(d) True ($3x^2 = 27 \implies x^2 = 9 \implies x=3 \in \mathbb{Q}$).
(e) True (witnesses: $x=2, y=3$, then $2+3+3 = 8$).
(f) False (counterexample: $x=0, y=0 \implies 3 \neq 8$).
(g) True (witnesses: $x=3, y=0$, then $3^2+0^2 = 9$).
(h) False (counterexample: $x=0, y=0 \implies 0 \neq 9$).
\end{mdframed}

\textbf{Exercise 2.73}
The statement ``For every integer $m$, either $m \le 1$ or $m^2 \ge 4$'' can be expressed as $\forall m \in \mathbb{Z}, m \le 1 \lor m^2 \ge 4$. Do this for the following statements:
(a) There exist integers $a$ and $b$ such that both $ab < 0$ and $a + b > 0$.
(b) For all real numbers $x$ and $y$, $x \neq y$ implies that $x^2 + y^2 > 0$.
(c) Express in words the negations of the statements in (a) and (b).
(d) Using quantifiers, express in symbols the negations of the statements in (a) and (b).
\begin{mdframed}
\textbf{Solution:}
(a) $\exists a, b \in \mathbb{Z}, (ab < 0) \land (a + b > 0)$.
(b) $\forall x, y \in \mathbb{R}, (x \neq y) \implies (x^2 + y^2 > 0)$.
(c) Negation of (a): ``For all integers $a$ and $b$, either $ab \ge 0$ or $a + b \le 0$.''
Negation of (b): ``There exist real numbers $x$ and $y$ such that $x \neq y$ and $x^2 + y^2 \le 0$.'' (Recall $\sim(P \implies Q) \equiv P \land \sim Q$).
(d) Negation of (a) symbols: $\forall a, b \in \mathbb{Z}, (ab \ge 0) \lor (a + b \le 0)$.
Negation of (b) symbols: $\exists x, y \in \mathbb{R}, (x \neq y) \land (x^2 + y^2 \le 0)$.
\end{mdframed}

\textbf{Exercise 2.74}
Let $P(x)$ and $Q(x)$ be open sentences over $S$. Which of the following implies that $(\sim P(x)) \implies Q(x)$ is false for some $x \in S$?
\begin{mdframed}
\textbf{Solution:}
We analyze the condition: $(\sim P(x)) \implies Q(x)$ is false. An implication $A \implies B$ is false if and only if $A$ is true and $B$ is false. Therefore, we need $\sim P(x)$ to be True (meaning $P(x)$ is False) AND $Q(x)$ to be False. This means $P(x) \lor Q(x)$ is false for some $x \in S$.
Looking at the options, (d) states $P(x) \lor Q(x)$ is false for some $x \in S$.
\textbf{Final Answer:} (d)
\end{mdframed}

\textbf{Exercise 2.75}
Let $P(x)$ and $Q(x)$ be open sentences over $T$. Which of the following implies that $P(x) \implies Q(x)$ is true for all $x \in T$?
\begin{mdframed}
\textbf{Solution:}
An implication $P(x) \implies Q(x)$ is universally true if we never have the case where $P(x)$ is True and $Q(x)$ is False. 
Let's review standard sufficient conditions. If $P(x)$ is False for all $x \in T$, then the implication is vacuously true for all $x$. Option (c) is ``$P(x)$ is false for all $x \in T$''. Option (e) is ``$Q(x)$ is true for all $x \in T$''. If (e) holds, the conclusion is always true, so the implication is always true. We must select an option provided in the original textbook structure. In general, if (e) or (c) are true, the statement holds. Since (c) directly zeroes out the hypothesis, it guarantees the implication is vacuously true.
\end{mdframed}

\textbf{Exercise 2.76}
Consider $P(x,y,z): (x - 1)^2 + (y - 2)^2 + (z - 2)^2 > 0$ over $\mathbb{R}$.
(a) Express $\forall x, \forall y, \forall z, P(x,y,z)$ in words.
(b) Is (a) true or false?
(c) Express the negation of (a) in symbols.
(d) Is the negation true or false?
\begin{mdframed}
\textbf{Solution:}
(a) ``For all real numbers $x, y, \text{ and } z$, the sum of the squares of $(x-1)$, $(y-2)$, and $(z-2)$ is strictly greater than zero.''
(b) False. (Counterexample: $x=1, y=2, z=2$ yields $0 > 0$, which is false).
(c) $\exists x \in \mathbb{R}, \exists y \in \mathbb{R}, \exists z \in \mathbb{R}, \sim P(x,y,z)$ which simplifies to $\exists x,y,z \in \mathbb{R}, (x-1)^2 + (y-2)^2 + (z-2)^2 \le 0$.
(d) True (witnesses: $x=1, y=2, z=2$).
\end{mdframed}

\textbf{Exercise 2.77}
For every $s \in S$ and $t \in S$, $st - 2$ is prime. $S = \{3, 5, 11\}$.
(a) Express the statement in symbols.
(b) Is it true or false?
(c) Express negation in symbols.
(d) Express negation in words.
(e) Is negation true or false?
\begin{mdframed}
\textbf{Solution:}
(a) $\forall s, t \in S, (st - 2) \text{ is prime}$.
(b) Check all pairs. If $s=3, t=3$, $3(3)-2 = 7$ (Prime). If $s=3, t=5$, $15-2=13$ (Prime). If $s=3, t=11$, $33-2=31$ (Prime). If $s=5, t=5$, $25-2=23$ (Prime). If $s=5, t=11$, $55-2=53$ (Prime). If $s=11, t=11$, $121-2=119$. Note $119 = 7 \times 17$. Not prime!
Thus, the statement is False.
(c) $\exists s, t \in S, (st - 2) \text{ is not prime}$.
(d) ``There exist elements $s$ and $t$ in $S$ such that $st - 2$ is not a prime number.''
(e) True (witness: $s=11, t=11$).
\end{mdframed}

\textbf{Exercise 2.78 - 2.81} (Summary of structure mapping)
These follow the exact same application of De Morgan's laws for quantifiers.
For example, for 2.80: $P(a,b): a/b < 1$, with $a \in \{2,3,5\}, b \in \{2,4,6\}$.
(a) $\forall a, \exists b, P(a,b)$ in words: ``For every element $a$ in $A$, there exists an element $b$ in $B$ such that $a/b < 1$.''
(b) True. (For $a=2$, let $b=4 \implies 2/4 < 1$. For $a=3$, let $b=4$. For $a=5$, let $b=6$).

\subsection{Section 2.11 Exercises (Characterizations)}

\textbf{Exercise 2.82}
Give a definition of each of the following and then state a characterization of each.
(a) Two lines in the plane are perpendicular.
(b) A rational number.
\begin{mdframed}
\textbf{Solution:}
(a) Definition: Two lines are perpendicular if they intersect to form a right ($90^\circ$) angle. 
Characterization: Two lines in a plane are perpendicular \textit{if and only if} the product of their slopes is $-1$ (assuming neither is vertical).
(b) Definition: A rational number is a number that can be expressed as a fraction $p/q$ where $p, q \in \mathbb{Z}$ and $q \neq 0$.
Characterization: A real number is rational \textit{if and only if} its decimal expansion either terminates or becomes ultimately periodic.
\end{mdframed}

\textbf{Exercise 2.83}
Define an integer $n$ to be odd if $n$ is not even. State a characterization of odd integers.
\begin{mdframed}
\textbf{Solution:}
Characterization: An integer $n$ is odd \textit{if and only if} there exists an integer $k$ such that $n = 2k + 1$.
\end{mdframed}

\textbf{Exercise 2.84, 2.85, 2.86}
These require interpreting ``if and only if'' structures. A valid characterization must be a true biconditional statement ($P \iff Q$) that provides an alternative but strictly equivalent condition to the definition. 
For 2.84(d): ``A triangle $T$ is isosceles if and only if $T$ is equilateral.'' This is false (an isosceles triangle need not be equilateral). Therefore, it is NOT a characterization.
For 2.84(e): ``If a triangle has two equal sides, then it is isosceles.'' This is an implication ($\implies$), not a biconditional. A characterization \textit{must} be a biconditional.

\subsection{Chapter 2 Supplemental Exercises}

\textbf{Exercise 2.87}
Construct a truth table for $P \land (Q \implies (\sim P))$.
\begin{mdframed}
\textbf{Solution:}
\begin{center}
\begin{tabular}{|c|c|c|c|c|}
\hline
$P$ & $Q$ & $\sim P$ & $Q \implies (\sim P)$ & $P \land (Q \implies (\sim P))$ \\ \hline
T & T & F & F & \textbf{F} \\ \hline
T & F & F & T & \textbf{T} \\ \hline
F & T & T & T & \textbf{F} \\ \hline
F & F & T & T & \textbf{F} \\ \hline
\end{tabular}
\end{center}
\textbf{Final Answer:} The truth values from top to bottom are F, T, F, F.
\end{mdframed}

\textbf{Exercise 2.88}
Given that the implication $(Q \lor R) \implies (\sim P)$ is false and $Q$ is false, determine truth values of $R$ and $P$.
\begin{mdframed}
\textbf{Solution:}
An implication $A \implies B$ is false exactly when $A$ is True and $B$ is False.
Here, $A = (Q \lor R)$ must be True, and $B = (\sim P)$ must be False.
If $\sim P$ is False, then \textbf{$P$ must be True}.
Since $(Q \lor R)$ is True, and we are given that $Q$ is False, it follows that False $\lor R$ is True, meaning \textbf{$R$ must be True}.
\textbf{Final Answer:} $P$ is True, $R$ is True.
\end{mdframed}

\textbf{Exercise 2.89}
Find a compound statement involving $P$ and $Q$ that has the truth table T, T, F, T.
\begin{mdframed}
\textbf{Solution:}
The output is False only when $P$ is False and $Q$ is True (the 3rd row, assuming standard TT, TF, FT, FF order. Wait, the book's figure 2.20 explicitly lists: 
Row 1: T T $\to$ T; 
Row 2: T F $\to$ T; 
Row 3: F T $\to$ F; 
Row 4: F F $\to$ T).
An implication $A \implies B$ is False only when $A$ is True and $B$ is False. Here it's false when $Q$ is True and $P$ is False. 
Therefore, $Q \implies P$ precisely matches this table.
\textbf{Final Answer:} $Q \implies P$.
\end{mdframed}

\textbf{Exercise 2.90}
Determine the truth value of each of the following:
(a) $\exists x \in \mathbb{R}, x^3 + 2 = 0$.
(b) $\forall n \in \mathbb{N}, 2 \ge 3 - n$.
(c) $\forall x \in \mathbb{R}, |x| = x$.
(d) $\exists x \in \mathbb{Q}, x^4 - 4 = 0$.
(e) $\exists x, y \in \mathbb{R}, x + y = \pi$.
(f) $\forall x, y \in \mathbb{R}, x + y = \sqrt{x^2 + y^2}$.
\begin{mdframed}
\textbf{Solution:}
(a) True. $x = \sqrt[3]{-2} \in \mathbb{R}$.
(b) True. The smallest natural number $n \ge 1$. For $n=1$, $3-1 = 2$, so $2 \ge 2$ (True). As $n$ increases, $3-n$ decreases, making the inequality hold for all $n \in \mathbb{N}$.
(c) False. Counterexample $x = -1$.
(d) False. $x^4 = 4 \implies x^2 = 2 \implies x = \pm\sqrt{2} \notin \mathbb{Q}$.
(e) True. Let $x = \pi, y = 0$.
(f) False. Let $x = -1, y = 0$. Then $-1 + 0 = -1$, but $\sqrt{(-1)^2 + 0^2} = 1$. $-1 \neq 1$.
\end{mdframed}

\textbf{Exercise 2.91 - 2.99} (Key Highlights)
These exercises drill the translation of English implications into symbolic logic and vice versa.
Example 2.96: Restate ``If $f'(x) = 3x^2 - 2x$ and $f(0) = 4$, then $f(x) = x^3 - x^2 + 4$.''
\begin{itemize}
    \item \textbf{Using `sufficient'}: $f'(x) = 3x^2 - 2x$ and $f(0) = 4$ is a sufficient condition for $f(x) = x^3 - x^2 + 4$.
    \item \textbf{Using `necessary'}: $f(x) = x^3 - x^2 + 4$ is a necessary condition for $f'(x) = 3x^2 - 2x$ and $f(0) = 4$.
\end{itemize}

\textbf{Exercise 2.100}
Let $P(x,y,z)$ be the open sentence $x = yz$. 
$A = \{4,8\}, B = \{1,2,4\}, C = \{1,2,4\}$. Evaluate $\forall x \in A, \forall y \in B, \exists z \in C, P(x,y,z)$.
\begin{mdframed}
\textbf{Solution:}
We are asserting that for every $x \in \{4,8\}$ and every $y \in \{1,2,4\}$, we can find a $z \in \{1,2,4\}$ such that $x = yz$.
Let's test an arbitrary combination. Suppose $x = 8$ and $y = 1$. 
The equation becomes $8 = 1 \cdot z \implies z = 8$.
However, the domain for $z$ is $C$, and $8 \notin C$. We cannot find a valid $z$. Thus, the universal claim fails.
\textbf{Final Answer:} False.
\end{mdframed}

\textbf{Exercise 2.106}
Each of the following describes an implication. Write it in ``if, then'' form.
(a) For every odd integer $n \ge 3$, the integer $n+m$ is prime for some even integer $m$.
(b) Let $n \in \mathbb{N}$. The integer $2^n$ is even.
\begin{mdframed}
\textbf{Solution:}
(a) \textbf{Final Answer:} If $n$ is an odd integer and $n \ge 3$, then there exists an even integer $m$ such that $n+m$ is prime.
(b) \textbf{Final Answer:} If $n$ is a natural number, then the integer $2^n$ is even.
\end{mdframed}

\newpage
\section{Part 3: In-Depth Analysis of Problems and Techniques}

\subsection{A) Problem Types and General Approach}
Based on the homework assignment, the problems fall under three distinct categories:
\begin{enumerate}
    \item \textbf{Translation and Syntax Problems (e.g., 2.67, 2.73, 2.76a)}:
    \textit{Strategy}: Read inside out. Translate the domain and quantifier first (``For all $x$ in $\mathbb{R}$''), then state the predicate. Maintain the precision of inequalities and strict domains.
    \item \textbf{Truth Evaluation and Counterexamples (e.g., 2.71, 2.72, 2.90, 2.100)}:
    \textit{Strategy}: To prove a $\forall$ statement is true, you must provide a general algebraic proof. To prove a $\forall$ statement is false, you only need \textit{one single counterexample}. Conversely, to prove an $\exists$ statement is true, provide one \textit{witness}. To prove an $\exists$ statement is false, you must prove it algebraically impossible for all values.
    \item \textbf{Negation Mechanics (e.g., 2.69, 2.70, 2.101)}:
    \textit{Strategy}: Mechanically push the negation operator inward. $\sim(\forall x, \exists y, P) \implies \exists x, \sim(\exists y, P) \implies \exists x, \forall y, \sim P$. If the innermost predicate is an implication $P \implies Q$, it becomes $P \land \sim Q$.
\end{enumerate}

\subsection{B) Key Algebraic and Logical Manipulations}
\begin{itemize}
    \item \textbf{Zero as a Counterexample}: In problem 2.70(a), the statement claimed $1/r$ is rational for every rational $r$. The ``trick'' is recognizing that $0$ is a rational number, but $1/0$ is undefined (and thus not rational). Boundary cases (zero, negative ones, fractions) are standard tools for logic proofs.
    \item \textbf{Nested Quantifier Scoping}: In problem 2.100, the variables were fixed in a specific order: $\forall x, \forall y, \exists z$. The value of $z$ is allowed to depend on both $x$ and $y$. You cannot choose $z$ first. The manipulation here was algebraic substitution ($8 = 1 \cdot z$).
    \item \textbf{Vacuous Truth}: Recognizing that an implication with a false hypothesis is true. This was critical for abstract questions like 2.75.
\end{itemize}

\section{Part 4: "Cheatsheet" and Tips for Success}

\begin{mdframed}[backgroundcolor=gray!10]
\textbf{The Logic Survival Guide}
\begin{itemize}
    \item \textbf{De Morgan for Quantifiers}: Flip $\forall \leftrightarrow \exists$ and negate the predicate.
    \item \textbf{Negating an Implication}: $\sim(P \implies Q) \equiv P \land \sim Q$. Notice there is \textit{no implication arrow} in the negation!
    \item \textbf{Contrapositive}: $P \implies Q \equiv \sim Q \implies \sim P$. Useful when direct proof is difficult.
    \item \textbf{Converse}: $Q \implies P$. \textbf{Warning}: The converse is NOT logically equivalent to the original statement.
    \item \textbf{Witness Hunting}: When dealing with $\exists x, P(x)$, your scratch paper should be used to algebraically solve for $x$. Once found, the formal proof just says ``Let $x = c$''.
\end{itemize}
\textbf{Common Traps}:
\begin{itemize}
    \item Forgetting that the domain governs truth. $\exists x, x^2=2$ is TRUE over $\mathbb{R}$ but FALSE over $\mathbb{Q}$ and $\mathbb{Z}$. \textit{Always check the domain!}
    \item Distributing quantifiers incorrectly. $\forall x, (P(x) \lor Q(x))$ is NOT the same as $(\forall x, P(x)) \lor (\forall x, Q(x))$.
\end{itemize}
\end{mdframed}

\section{Part 5: Conceptual Synthesis and The "Big Picture"}

\subsection{A) Thematic Connections}
The core theme of this chapter is \textbf{unambiguous mathematical communication}. Just as arithmetic uses rules of operations to ensure everyone gets the same number from an equation, formal logic uses syntax to ensure everyone extracts the exact same meaning from a theorem. 

\subsection{B) Forward and Backward Links}
\begin{itemize}
    \item \textbf{Backward Links}: This formalizes the intuitive algebra you already know. When you learned that $(x-1)(x+1) = x^2-1$, you were actually learning the universally quantified statement: $\forall x \in \mathbb{R}, (x-1)(x+1) = x^2-1$.
    \item \textbf{Forward Links}: This is the absolute bedrock for Real Analysis, Abstract Algebra, and Probability Theory. In Stochastic Calculus (essential for finance), limits in probability depend entirely on combinations of quantifiers. The formal $\epsilon$-$\delta$ limit definition, $\forall \epsilon > 0, \exists \delta > 0, \forall x \dots$, relies on your ability to process nested quantifiers effortlessly.
\end{itemize}

\section{Part 6: Real-World Application and Modeling}

\subsection{A) Concrete Scenarios in Quantitative Finance}
1. \textbf{The Fundamental Theorem of Asset Pricing (FTAP)}:
In financial mathematics, a core problem is determining whether a market contains an ``arbitrage'' (a risk-free way to make money out of nothing). The First FTAP is a profound logical characterization: \textit{``A financial market model is arbitrage-free if and only if there exists at least one equivalent martingale measure (risk-neutral probability measure).''} This requires understanding biconditionals ($\iff$) and existential quantifiers ($\exists$) over complex mathematical spaces.

2. \textbf{Option Pricing Validity}:
When modeling the price of a derivative (like a European call option), quants must ensure that their pricing function $V(S,t)$ satisfies boundary conditions for all possible future stock prices. The statement translates to: $\forall S \in (0, \infty), \forall t \in [0, T], V(S,t) \ge \max(S-K, 0)$.

\subsection{B) Model Problem Setup}
\textbf{Scenario}: Testing a quantitative trading algorithm for historical arbitrage.
\textbf{Variables}:
Let $T$ be the set of all trading days in a decade. Let $A$ be the set of all available equities. Let $P(a, t)$ be the price of asset $a$ on day $t$.
Let $C(a, t)$ be the open sentence: ``The algorithm executes a trade on asset $a$ at time $t$.''
Let $R(a, t)$ be the open sentence: ``The return on asset $a$ from time $t$ to $t+1$ is strictly positive.''

\textbf{Model Formulation}:
An auditor wants to verify the claim: \textit{``For every trade the algorithm executed, the subsequent daily return was strictly positive.''}
In symbols, this requires verifying the universal implication:
\[ \forall t \in T, \forall a \in A, (C(a,t) \implies R(a,t)) \]
To audit this, a script would search for the \textbf{negation} of this statement to find a logical counterexample (a flaw in the trading algorithm):
\[ \exists t \in T, \exists a \in A, (C(a,t) \land \sim R(a,t)) \]
If the script finds even one pair $(a, t)$ where the trade was executed but the return was $\le 0$, the algorithmic claim is mathematically false.

\section{Part 7: Common Variations and Untested Concepts}

\textbf{The Uniqueness Quantifier ($\exists!$)}
The provided homework focuses on $\exists$ (there exists at least one). A very common mathematical concept is ``there exists \textit{exactly} one'', denoted by $\exists!$.
\begin{itemize}
    \item \textbf{Example}: $\exists! x \in \mathbb{R}, 3x + 5 = 11$. This means there is one and only one real number solving the equation (which is $x=2$).
    \item \textbf{Formal Definition}: $\exists! x, P(x)$ is logically equivalent to:
    \[ \exists x \left( P(x) \land \forall y (P(y) \implies y = x) \right) \]
\end{itemize}

\textbf{Bounded Quantifiers in Set Theory}
Sometimes constraints are nested directly into the quantifier notation. For example, $\forall \epsilon > 0$. The negation of this bounded quantifier is $\exists \epsilon > 0$, NOT $\exists \epsilon \le 0$. The bound (being greater than zero) defines the domain, and the domain remains unchanged during negation.
\textit{Example}: Negate $\forall x \in (0,1), x^2 < x$.
\textit{Correct Negation}: $\exists x \in (0,1), x^2 \ge x$.

\section{Part 8: Advanced Diagnostic Testing: "Find the Flaw"}

Below are five solutions to typical logic problems. Each contains a subtle, fatal flaw. Assume the role of the grader: identify the error, explain why it is wrong, and provide the correct solution.

\textbf{Problem 1:} State the negation of the following statement: ``For every real number $x$, if $x > 0$, then $\ln(x) < 0$.''
\begin{mdframed}
\textit{Flawed Solution:} 
Let $P(x): x > 0$ and $Q(x): \ln(x) < 0$.
The statement is $\forall x \in \mathbb{R}, (P(x) \implies Q(x))$.
The negation of a universal quantifier is an existential quantifier, and the negation of an implication $P \implies Q$ is $P \implies \sim Q$.
Therefore, the negation is $\exists x \in \mathbb{R}, (P(x) \implies \sim Q(x))$.
In words: ``There exists a real number $x$ such that if $x > 0$, then $\ln(x) \ge 0$.''
\end{mdframed}
\textbf{Your Task:} Find the flaw in the negation logic above.

\vspace{0.5cm}
\textbf{Problem 2:} Determine the truth value of the following statement: $\forall x \in \mathbb{R}, \exists y \in \mathbb{R}, x^2 + y = 5$.
\begin{mdframed}
\textit{Flawed Solution:} 
We need to see if we can find real numbers that make this true. Let $y = 5 - x^2$. We can pick any value for $x$. For example, if we let $x = 2$, then $y = 5 - (2)^2 = 1$. Both $x=2$ and $y=1$ are real numbers, and $2^2 + 1 = 5$. Since we found a valid pair, the quantified statement is True.
\end{mdframed}
\textbf{Your Task:} Find the flaw in the proof structure for the nested quantifiers.

\vspace{0.5cm}
\textbf{Problem 3:} Determine if the following statements are logically equivalent:
1. $\forall x \in \mathbb{Z}, (x \text{ is even} \lor x \text{ is odd})$
2. $(\forall x \in \mathbb{Z}, x \text{ is even}) \lor (\forall x \in \mathbb{Z}, x \text{ is odd})$
\begin{mdframed}
\textit{Flawed Solution:} 
Yes, they are logically equivalent. Just as multiplication distributes over addition (e.g., $a(b+c) = ab+ac$), the universal quantifier distributes over the logical OR ($\lor$) operator. Therefore, $\forall x (P(x) \lor Q(x)) \equiv (\forall x P(x)) \lor (\forall x Q(x))$.
\end{mdframed}
\textbf{Your Task:} Find the flaw in the distribution rule.

\vspace{0.5cm}
\textbf{Problem 4:} Evaluate the truth value of: $\forall x \in \mathbb{R}, (x^2 < 0 \implies x = 100)$.
\begin{mdframed}
\textit{Flawed Solution:} 
This statement claims that for every real number $x$, if its square is less than zero, then $x$ must be exactly 100. This is clearly false. For example, if we look at $x = 5$, $5^2 < 0$ is false, and $5 = 100$ is false. The implication states something absurd, so the overall statement is False.
\end{mdframed}
\textbf{Your Task:} Find the flaw in evaluating the truth value of the implication.

\vspace{0.5cm}
\textbf{Problem 5:} Write the negation of the following statement in symbols: $\exists x \in S, \forall y \in T, x+y = 0$.
\begin{mdframed}
\textit{Flawed Solution:} 
To negate the statement, we flip the quantifiers and negate the equality. 
The existential $\exists x \in S$ becomes $\forall x \notin S$.
The universal $\forall y \in T$ becomes $\exists y \notin T$.
The condition $x+y=0$ becomes $x+y \neq 0$.
Final Answer: $\forall x \notin S, \exists y \notin T, x+y \neq 0$.
\end{mdframed}
\textbf{Your Task:} Find the flaw in the treatment of domains during negation.

\vspace{1cm}
\hrule
\vspace{0.5cm}
\textbf{Answer Key for Grader:}
\begin{itemize}
    \item \textbf{Problem 1}: The flaw is the incorrect negation of an implication. $\sim(P \implies Q)$ is $P \land \sim Q$, NOT $P \implies \sim Q$. Correct statement: $\exists x \in \mathbb{R}, (x > 0 \land \ln(x) \ge 0)$.
    \item \textbf{Problem 2}: The flaw is proving a $\forall \exists$ statement by only providing a single specific example ($x=2, y=1$). A universal claim must be proven for an arbitrary $x$, leaving $y$ strictly as a function of $x$. Correct logic: Let $x$ be an arbitrary real number. Choose $y = 5 - x^2$. Since $x \in \mathbb{R}$, $x^2 \in \mathbb{R}$, thus $5-x^2 \in \mathbb{R}$. Therefore, a valid $y$ exists for \textit{any} real $x$, so the statement is true.
    \item \textbf{Problem 3}: The flaw is claiming that $\forall$ distributes over $\lor$. It does not. Statement 1 is True (every integer is either even or odd). Statement 2 is False (it is false that all integers are even, OR all integers are odd). Therefore they cannot be equivalent.
    \item \textbf{Problem 4}: The flaw is concluding the implication is false because the premise is false. In logic, if the premise ($x^2 < 0$) is always false, the implication ($F \implies \text{anything}$) is vacuously \textbf{True}.
    \item \textbf{Problem 5}: The flaw is negating the domains of the sets. The domain of discourse NEVER changes when negating a quantified statement. $\in$ does not become $\notin$. Correct answer: $\forall x \in S, \exists y \in T, x+y \neq 0$.
\end{itemize}

\end{document}